\documentclass[11pt]{article}
\usepackage[spanish,es-nodecimaldot]{babel}
\usepackage[utf8]{inputenc}
\usepackage{amsmath,amssymb,geometry,array,tikz}
\geometry{letterpaper,margin=1.5cm}
\AtBeginDocument{\setlength{\abovedisplayskip}{3pt}\setlength{\belowdisplayskip}{3pt}\setlength{\abovedisplayshortskip}{2pt}\setlength{\belowdisplayshortskip}{2pt}}
\setlength{\parindent}{0pt}
\renewcommand{\arraystretch}{1.3}
\DeclareMathOperator{\sen}{sen}
\begin{document}
\begin{center}{\Large\textbf{Simulacro 01 --- Matemáticas Aplicadas}}\end{center}
\hrule\vspace{8pt}

%=================== EJERCICIO 1 ===================
{\large\textbf{1) Solución de E.D.O de 2º con valor inicial usando R-K de 4º}}\\[6pt]
\textbf{Resolver} \quad $y''-3y'=4y+2\sen(t) \qquad y(1)=2 \quad y'(1)=3 \qquad H=0.1$\\[4pt]
\[
\begin{array}{r} f:\\[2pt] g:\end{array}\left\{\begin{array}{l} y'=z\\[2pt] z'=3z+4y+2\sen(t)\end{array}\right.
\qquad\qquad y(1)=2 \qquad z(1)=3
\]
\[
\begin{array}{ll}
K_1=Hf(t,y,z) & C_1=Hg(\,''\,)\\
K_2=Hf\!\left(t+\frac{H}{2},\,y+\frac{K_1}{2},\,z+\frac{C_1}{2}\right) & C_2=Hg(\,''\,)\\
K_3=Hf\!\left(t+\frac{H}{2},\,y+\frac{K_2}{2},\,z+\frac{C_2}{2}\right) & C_3=Hg(\,''\,)\\
K_4=Hf(t+H,\,y+K_3,\,z+C_3) & C_4=Hg(\,''\,)\\[4pt]
\multicolumn{2}{l}{y_{n+1}=y_n+\big(K_1+2(K_2+K_3)+K_4\big)/6 \qquad z_{n+1}=z_n+\big(C_1+2(C_2+C_3)+C_4\big)/6}
\end{array}
\]
{\small
%%EJ1%%
}
\vspace{-6pt}\[ %%RES1%% \]

\newpage
%=================== EJERCICIO 2 ===================
{\large\textbf{2) Solución de E.D.O de 2º con valor inicial usando diferencias finitas}}\\[6pt]
\textbf{Resolver} \quad $y''-6y'+10y=t \qquad y(1)=1 \quad y'(1)=0 \qquad n=5$\\[4pt]
$h=0.1$ \hfill
\begin{tikzpicture}[baseline=-3pt,x=1.25cm]
\draw[dashed] (-1,0)--(0,0); \draw (0,0)--(5,0);
\foreach \t/\i in {0/1,1/1.1,2/1.2,3/1.3,4/1.4,5/1.5}{\draw (\t,0.08)--(\t,-0.08); \node[above,font=\small] at (\t,0.08){\i};}
\foreach \k in {0,...,5}{\node[below,font=\small] at (\k,-0.08){$y_{\k}$};}
\draw (-1,0.08)--(-1,-0.08); \node[below,font=\small] at (-1,-0.08){$y_{-1}$}; \node[above,font=\scriptsize] at (-1,0.08){ficticio};
\end{tikzpicture}
\begin{align*}
&\frac{(y_{i+1}-2y_i+y_{i-1})}{h^2}-6\,\frac{(y_{i+1}-y_{i-1})}{2h}+10y_i=t_i\\[4pt]
&y_{i+1}-2y_i+y_{i-1}-3h(y_{i+1}-y_{i-1})+10h^2y_i=t_ih^2 && (A)\\[4pt]
&(1+3h)\,y_{i-1}-(2-10h^2)\,y_i+(1-3h)\,y_{i+1}=t_ih^2 && (B)
\end{align*}
$y'(1)=\dfrac{y_1-y_{-1}}{2h}=0 \;\Rightarrow\; y_{-1}=y_1$
\begin{align*}
i=0 \;\; t_0=1 &\quad (1+3(0.1))\,y_{-1}-(2-10(0.1)^2)\,y_0+(1-3(0.1))\,y_1=1(0.1)^2 && (0)\\
&\quad 1.3\,y_1-1.9(1)+0.7\,y_1=0.01 \;\to\; 2y_1=1.91\\
i=1 \;\; t_1=1.1 &\quad (1+3(0.1))\,y_0-(2-10(0.1)^2)\,y_1+(1-3(0.1))\,y_2=1.1(0.1)^2 && (1)\\
&\quad 1.3(1)-1.9\,y_1+0.7\,y_2=0.011 \;\to\; -1.9\,y_1+0.7\,y_2=-1.289\\
i=2 \;\; t_2=1.2 &\quad (1+3(0.1))\,y_1-(2-10(0.1)^2)\,y_2+(1-3(0.1))\,y_3=1.2(0.1)^2 && (2)\\
i=3 \;\; t_3=1.3 &\quad (1+3(0.1))\,y_2-(2-10(0.1)^2)\,y_3+(1-3(0.1))\,y_4=1.3(0.1)^2 && (3)\\
i=4 \;\; t_4=1.4 &\quad (1+3(0.1))\,y_3-(2-10(0.1)^2)\,y_4+(1-3(0.1))\,y_5=1.4(0.1)^2 && (4)
\end{align*}
\[
\underbrace{\begin{bmatrix}
2 & 0 & 0 & 0 & 0\\
-1.9 & 0.7 & 0 & 0 & 0\\
1.3 & -1.9 & 0.7 & 0 & 0\\
0 & 1.3 & -1.9 & 0.7 & 0\\
0 & 0 & 1.3 & -1.9 & 0.7
\end{bmatrix}}_{A}
\underbrace{\begin{bmatrix} y_1\\ y_2\\ y_3\\ y_4\\ y_5\end{bmatrix}}_{X}
=
\underbrace{\begin{bmatrix} 1.91\\ -1.289\\ 0.012\\ 0.013\\ 0.014\end{bmatrix}}_{B}
\qquad X=A^{-1}B
\]
\[
\boxed{\begin{array}{c|cccccc}
t & 1 & 1.1 & 1.2 & 1.3 & 1.4 & 1.5\\\hline
y & 1 & %%Y1%% & %%Y2%% & %%Y3%% & %%Y4%% & %%Y5%%
\end{array}}
\]
{\footnotesize $^{*}$ El enunciado no indica el intervalo; se toma $h=0.1$ (5 subintervalos, $1\le t\le1.5$).}

\newpage
%=================== EJERCICIO 3 ===================
{\large\textbf{3) Teoría de la variable compleja}}\\[6pt]
\textbf{Calcule} \quad $\displaystyle\int_0^{4+2i}(z-iz)^2\,dz$ \quad siguiendo las trayectorias $C_1,\,C_2,\,C_3$\\[6pt]
$z=x+iy$ \hfill
\begin{tikzpicture}[baseline=1cm,scale=0.8]
\draw[->] (-0.3,0)--(5,0) node[right]{$x$}; \draw[->] (0,-0.3)--(0,2.8) node[above]{$i$};
\draw[thick,->] (0,0)--(0,1.1); \draw[thick] (0,1)--(0,2); \node[left] at (0,1){$C_1$};
\draw[thick,->] (0,2)--(2.1,2); \draw[thick] (2,2)--(4,2); \node[above] at (2,2){$C_2$};
\draw[thick,->] (0,0)--(2.1,1.05); \draw[thick] (2,1)--(4,2); \node[below right] at (2,1){$C_3$};
\draw[dashed] (4,0)--(4,2);
\node[left] at (0,2){$(0,2)$}; \node[right] at (4,2){$(4,2)$}; \node[below left] at (0,0){$(0,0)$};
\end{tikzpicture}\\[4pt]
\textbf{Tray $C_1$:} \quad $x=0 \quad z=iy \;\to\; dz=i\,dy$
\begin{align*}
\int (z-iz)^2dz&=i\int_0^2(iy-i(iy))^2dy=i\int_0^2(y+iy)^2dy=i\int_0^2(y^2+2y^2i-y^2)\,dy\\
&=i\int_0^2 2y^2i\,dy=-\frac{2y^3}{3}\Big|_0^2=-\frac{16}{3}
\end{align*}
\textbf{Tray $C_2$:} \quad $y=2 \quad z=x+2i \;\to\; dz=dx$
\begin{align*}
\int_0^4\big(x+2i-i(x+2i)\big)^2dx&=\int_0^4\big((x+2)+(2-x)i\big)^2dx=\int_0^4\big((x+2)^2-(2-x)^2+2(x+2)(2-x)i\big)dx\\
&=\int_0^4\big(8x+8i-2x^2i\big)dx=4x^2+8xi-\frac{2x^3}{3}i\,\Big|_0^4=64+32i-\frac{128}{3}i=64-\frac{32}{3}i
\end{align*}
\[
C_1+C_2=-\frac{16}{3}+64-\frac{32}{3}i=\frac{176}{3}-\frac{32}{3}i
\]
\textbf{Tray $C_3$:} \quad $y=\dfrac{x}{2} \quad z=x+\dfrac{x}{2}i \;\to\; dz=\left(1+\dfrac{i}{2}\right)dx$
\begin{align*}
\int_0^4\left(x+\frac{x}{2}i-i\left(x+\frac{x}{2}i\right)\right)^2\left(1+\frac{i}{2}\right)dx&=\int_0^4\left(\frac{3x}{2}-\frac{x}{2}i\right)^2\left(1+\frac{i}{2}\right)dx=\int_0^4\left(2x^2-\frac{3}{2}x^2i\right)\left(1+\frac{i}{2}\right)dx\\
&=\int_0^4\left(\frac{11}{4}x^2-\frac{1}{2}x^2i\right)dx=\frac{11x^3}{12}-\frac{x^3}{6}i\,\Big|_0^4=\frac{176}{3}-\frac{32}{3}i
\end{align*}
\[
\boxed{\int_0^{4+2i}(z-iz)^2\,dz=\frac{176}{3}-\frac{32}{3}i\approx 58.6666-10.6666\,i \quad\text{(por } C_1+C_2 \text{ y por } C_3\text{)}}
\]
{\footnotesize $^{*}$ $(z-iz)^2=-2iz^2$ es analítica: la integral no depende de la trayectoria. Se toma $C_3$: $y=x/2$ (el texto ``$y=x^2+1$'' no pasa por $(0,0)$ ni por $(4,2)$).}
\end{document}
